\section{Scales and the local experiment}
\label{scale:section}

We now specify the scales, geometric randomizations, and common inputs
for the local updates. The next sections construct those updates and
prove the separation and persistence estimates stated above.

\subsection{Fixed constants and scales}\label{sets:parameters}
The geometric constants satisfy
$\epsilon\ll j\ll g\ll R\ll D\ll1$. The table records their roles;
the explicit choices below supply every numerical inequality used later.

\begin{center}
\renewcommand{\arraystretch}{1.12}
\begin{tabular}{@{}p{0.19\textwidth}p{0.77\textwidth}@{}}
\toprule
Parameters & Role \\
\midrule
$D$, $R$, $g$ & Main radial cell width, secondary update radius, and radial width for grouping seeds. \\
$j$, $\epsilon$ & Detectable drop at a separating tip, and error in a descending path. \\
$a_1$, $a_2$ & Broad and narrow band widths; control vertical transfers and interpolation. \\
$E_A$, $E_B$ & Interpolation errors in the two systems. \\
$\kappa_A$, $\kappa_B$ & Convert field variation into label mass while preserving vertical nesting. \\
$\lambda$ & Size of the auxiliary correction in system $B$ (Section~\ref{sec:increment}). \\
$Q$ & Scale ratio; controls earlier chart splitting and later movement. \\
\bottomrule
\end{tabular}
\end{center}

Fix the following numerical constants:
\[
 R=10^{-8},\quad D=10^{-3},\quad g=10^{-14},\quad
 j=10^{-20},\quad \epsilon=10^{-24},\quad
 E_A=\epsilon/1000,\quad E_B=10^{-3}R.
\]
Set $\zeta_A=\epsilon/100$ and $\zeta_B=1/10$. Choose $a_1>0$ within the
width bound of Theorem~\ref{int:interpolation} for $(\zeta_A,E_A)$, within the bounds of Lemmas~\ref{geo:cells} and~\ref{geo:padding} for $D$ and $g$, and with $a_1\le10^{-6}E_A$.
Set $\kappa_A=a_1/10^6$. Choose $a_2>0$ within the bound of Theorem~\ref{int:interpolation} for
$(\zeta_B,E_B)$ and with
\[
 a_2\le10^{-6}\min\{a_1,\kappa_Aj\},
 \qquad \kappa_B=a_2/10^6.
\]
Reserve $K=100$ as a Lipschitz bound for the auxiliary field in system
$B$; Lemma~\ref{sec:increment-bound} will supply this bound. Set
$\lambda=R/(10^{12}K)$. If necessary decrease $a_1,a_2$ further to meet the
explicit thin-band hypotheses of the geometric lemmas, and then choose
their dependent constants in the order just specified. All inequalities
in this choice are upper bounds on a newly chosen positive number.
The scale spacing $Q=2^m$ is chosen last.

Choose finitely many dyadic scales so that every nonzero vertex distance
belongs to $[r,2r)$. For each residue class of exponents modulo $m$, run one
schedule in decreasing order of its scales. For the analysis of a step,
divide distances, heights, fields, and their additive constants by its
scale $r$. Multiply the resulting threshold and field changes by $r$
before evaluating them in the original label space. In particular a field
change $\Delta P$ in scale units moves label mass
$r\kappa|\Delta P|$. The label coordinates are not resampled or
renormalized.

\subsection{The data of a step}\label{sets:setup}
At a step tile the height line by intervals $J_1$ of width $a_1$ with a
common random shift, and independently by intervals $J_2$ of width $a_2$.
For the two clipping constructions, process every broad or narrow band
meeting the height interval
\[
 [\min_{v\in V}h(v)-6,\ \max_{v\in V}h(v)+6]
\]
in the current scale units, and make no clipping update in other bands. This is a
finite collection and contains every band needed for a radius-$5$
neighborhood of an original vertex. The decision whether to process a
band is common to that entire band. This convention does not truncate the
ordered increments of Section~\ref{sec:increment}, which act on their
current chart regions. Given the bands, fix the radial basepoints in
their components before sampling any radial layer origins. In every
broad-band component form both the width-$D$ radial cells
and the width-$g$ radial cells, using independently shifted origins.
The narrow bands will instead determine the domains of the secondary
proposals and the balls in the packing estimate. Interpolation
basepoints and portals are likewise fixed before
any interpolation noises.

A vertex is \emph{interior} if its narrow band is contained in its broad
band, its height avoids the outer tenths of both bands, and all its column
representatives have radial values farther than $0.02D$ and $0.02g$ from
the respective layer interfaces. A pair with height difference less than
$10^{-4}a_2$ has probability at least $1/4$ of satisfying these requirements
in common bands. Indeed each of the two height tests has failure probability
at most $0.21$, and each radial test, for two vertices together, has failure
probability at most $0.081$, with the remaining allowance covering the
column-width and grid errors. A chosen representative controls all the
others within an error at most $a_1$, without a union bound over columns.
Interior in the broad band implies containment of the narrow band because
$a_2\le10^{-6}a_1$.

Recall that eligibility at the current scale requires the two closed
$d$-balls of radius $5$ to lie in one common old region in each system.
At this scale use the complete finite list
of eligible original vertex pairs with $1\le d(x,y)<2$. For the detailed
tests impose
\begin{equation}\label{sets:cutoffs}
 |h(x)-h(y)|<10^{-4}a_2,\quad
 |s_A(x)-s_A(y)|<E_A,\quad
 |P_B(x)-P_B(y)|<10^{-4}R.
\end{equation}
Here $s_A$ is the ordinary orientation and $P_B$ is the full old field.
Failure of a cutoff already gives positive separation by
\eqref{sets:distance}, so it is handled by the no-change mode.

There are four equally likely modes: no change, the main update in system
$A$, the secondary clipping in system $B$, and the ordered increment in
system $B$. The surviving groups and admission tests will be defined in
Sections~\ref{pack:seed-definition-section} and~\ref{sec:secondary}.
These lists, tests, and geometric data use the old charts; they are
determined independently of the choice of mode and the
fresh interpolation noises. Cell and group coins are independent fair
coins. Where the secondary construction in Section~\ref{sec:secondary} calls
for a fair sign, all points with one old chart identifier use the same
sign. These construction signs are independent of the label-coordinate
sign $\sigma$.

Thus a step first determines its bands and cells, then its eligible pairs,
surviving groups and admissions, and finally its fresh interpolation and
branch offsets. Conditioning on the earlier choices leaves those last
offsets independent. In particular, no admission decision is made by
examining a newly selected interpolation atom.

To avoid selection-measurability issues, all bounded-interval geometric
randomizations may use finite equally spaced grids. For a fixed finite
set of scales take their relative meshes at most $\eta_0$, where $\eta_0$
is smaller than every finer-to-coarser scale ratio and small enough for
the fixed interior-probability allowances above. An interval probability
bound $C\rho$ becomes $C'(\rho+\eta_0)$, which has the required size for
all comparisons in the proof. There are finitely many outcomes at every
step and finitely many nodes in each interpolation hierarchy. Choices
depending on past outcomes can therefore be made separately on each branch
of a finite tree. This makes every final random mixture a finite direct sum.
